\documentclass[10pt,a4paper]{amsart}
\usepackage[margin=19mm]{geometry}
\usepackage{amsmath,amssymb,mathtools}
\usepackage[scr=rsfs,cal=euler]{mathalfa}
\usepackage{xfrac}
\usepackage[unicode,hidelinks,bookmarks=false]{hyperref}
\newcommand{\ie}{i.e.\ }
\newcommand{\cf}{cf.\ }
\newcommand{\resp}{respectively\ }
\newcommand{\Subsection}{\subsection}
\providecommand{\nobreakdash}{\nobreak\hbox{-}}
\setlength{\emergencystretch}{2em}
\multlinegap0pt
\usepackage[shortlabels]{enumitem}
\usepackage{mathtools}
\usepackage{amssymb}
\usepackage{xcolor}
\usepackage{algorithm}\usepackage{algpseudocode}
\usepackage{tikz}
\newtheorem{theorem}{Theorem}
\title{Mode-Wise Spectral Criteria for Coupled Mass Transport in Hybrid PDE–ODE Tumor Microenvironments}
\author{Jiguang Yu, Louis Shuo Wang, Zonghao Liu, Jingfeng Liu}
\date{Source: arXiv:2601.20204v1 (CC BY 4.0)}
\begin{document}
\maketitle

\noindent\emph{Source and licence.} Mode-Wise Spectral Criteria for Coupled Mass Transport in Hybrid PDE–ODE Tumor Microenvironments. Version arXiv:2601.20204v1 (CC BY 4.0), distributed by arXiv under the Creative Commons Attribution 4.0 International licence, http://creativecommons.org/licenses/by/4.0/, as recorded in the arXiv licence metadata for this exact version and shown on the abstract page https://arxiv.org/abs/2601.20204v1. Full licence notice: \texttt{licenses/arxiv-2601.20204-CC-BY-4.0.txt}.\par\smallskip

\noindent\textit{Editorial scope.} Counted unit: the article's theorem thm:twoway\_criteria (source0.tex lines 836--907). The article's complete original proof of it is delivered with it (source0.tex lines 1335--1373, one \texttt{proof} environment of 39 lines, which the article places away from the statement). No mathematics is written, repaired or re-derived: the statement and the proof are the article's own lines, in the article's own notation, and no retained line is substituted or re-flowed.  Retained uncounted as the source-local context the proof needs: lines 824--831.  Nothing was omitted from any retained span, so the package carries no omission record.  No retained line contains a citation, so the wrapper carries no bibliography and no bibliography entry.  No retained line contains a figure environment, an image-inclusion command or drawing code, so no image is delivered and the package declares no asset dependency.  Editorial formatting only, in addition: an AMS \texttt{amsart} preamble that restates the article's own declarations and macros used by the retained lines, namely the environments \texttt{theorem} on separate counters, together with the standard packages those lines need.  The retained environments print in this document as Theorem 1 in order of appearance, whereas the article prints the same items with the numbering of its own full document.  The complete native source of the article as distributed in the arXiv e-print for this exact version is bundled unchanged as \texttt{sources/193511/source0.tex}.
\begin{equation}
\label{eq:twoway_reduced_g}
\begin{cases}
\partial_t S = d_S\Delta S - \chi_S\nabla\cdot\bigl(S\nabla g(S,R)\bigr) + f_S(S,R),\\[3pt]
\partial_t R = d_R\Delta R - \chi_R\nabla\cdot\bigl(R\nabla g(S,R)\bigr) + f_R(S,R),
\end{cases}
\qquad \partial_{\mathrm n}S=\partial_{\mathrm n}R=0,
\end{equation}

\begin{theorem}[Two-way feedback recovers Turing-type instability: trace/determinant criteria]
\label{thm:twoway_criteria}
Assume \eqref{eq:twoway_reduced_g} admits a spatially homogeneous interior equilibrium $\mathbf{u}^*=(S^*,R^*)$.
Let
\[
\mathcal{J}_F \coloneqq D(f_S,f_R)(S^*,R^*)=
\begin{pmatrix}
a & b\\
c & d
\end{pmatrix},
\qquad
g_S^*\coloneqq \partial_S g(S^*,R^*),\quad g_R^*\coloneqq \partial_R g(S^*,R^*),
\]
and assume kinetic stability:
\begin{equation}
\label{eq:kinetic_stability_twoway}
\operatorname{tr}(\mathcal{J}_F)<0,\qquad \det(\mathcal{J}_F)>0.
\end{equation}
For each Neumann mode $k\ge2$ with eigenvalue $\lambda_k>0$, the mode matrix is
\[
\mathcal{A}(\lambda_k)=\mathcal{J}_F-\lambda_k(\mathcal{D}-\mathcal{H}),
\qquad
\mathcal{D}\coloneqq \mathrm{diag}(d_S,d_R),
\]
where
\[
\mathcal{H}\coloneqq
\begin{pmatrix}
\chi_S g_S^* S^* & \chi_S g_R^* S^*\\[4pt]
\chi_R g_S^* R^* & \chi_R g_R^* R^*
\end{pmatrix}.
\]
Then the equilibrium $\mathbf{u}^*$ is unstable for some nonconstant mode (i.e.\ a Turing-type instability occurs) if and only if
there exists $k\ge2$ such that at least one of the following holds:
\begin{enumerate}[label=(\alph*), leftmargin=*]
\item \textbf{Trace condition:}
\begin{equation}
\label{eq:twoway_trace_condition}
\operatorname{tr}(\mathcal{A}(\lambda_k))
=
\operatorname{tr}(\mathcal{J}_F)
-\lambda_k(d_S+d_R)
+\lambda_k\bigl(\chi_S g_S^* S^*+\chi_R g_R^* R^*\bigr)
>0.
\end{equation}

\item \textbf{Determinant condition:}
$\det(\mathcal{A}(\lambda_k))<0$ for some $k\ge2$.
Equivalently, writing
\begin{equation}
\label{eq:det_quadratic_def}
\det(\mathcal{A}(\lambda))=A_1\lambda^2+A_2\lambda+\det(\mathcal{J}_F),
\qquad \lambda\ge0,
\end{equation}
with
\begin{align}
\label{eq:A1A2_def}
A_{1}
&= d_{S}d_{R}-d_{S}\chi_{R}g_R^*R^*-d_{R}\chi_{S}g_S^*S^*,\nonumber\\
A_{2}
&= -ad_{R}-dd_{S}
+a\chi_{R}g_R^*R^* + d\chi_{S}g_S^*S^*
-b\chi_{R}g_S^*R^* -c\chi_{S}g_R^*S^*,
\end{align}
there exists $k\ge2$ with $\det(\mathcal{A}(\lambda_k))<0$ if and only if one of the following (mutually exclusive) cases holds:
\begin{enumerate}[label=(\roman*), leftmargin=*]
\item $A_1<0$ (then $\det(\mathcal{A}(\lambda))\to -\infty$ as $\lambda\to\infty$ and large modes destabilize);
\item $A_1=0$ and $A_2<0$ (then $\det(\mathcal{A}(\lambda))=A_2\lambda+\det(\mathcal{J}_F)$ becomes negative for sufficiently large $\lambda$);
\item $A_1>0$, $A_2<0$, and $A_2^2-4A_1\det(\mathcal{J}_F)>0$ (then the upward-opening quadratic attains negative values between its two positive roots).
\end{enumerate}
\end{enumerate}
\end{theorem}

\begin{proof}
Consider the reduced two-way closure \eqref{eq:twoway_reduced_g} and let $\mathbf{u}^*=(S^*,R^*)$ be a homogeneous interior equilibrium.
Linearizing at $\mathbf{u}^*$ yields
\[
\partial_t \tilde{\mathbf{u}}=(\mathcal{D}-\mathcal{H})\Delta \tilde{\mathbf{u}}+\mathcal{J}_F\tilde{\mathbf{u}},
\qquad
\tilde{\mathbf{u}}=(S-S^*,R-R^*)^{\mathsf T},
\]
where $\mathcal{D}=\mathrm{diag}(d_S,d_R)$, $\mathcal{J}_F=DF(\mathbf{u}^*)$, and $\mathcal{H}$ is given in Theorem~\ref{thm:twoway_criteria}.

Projecting onto the Neumann eigenmode $\omega_k$ with $\lambda_k>0$ gives the modal ODE
\[
\frac{d}{dt}\hat{\mathbf{u}}=\mathcal{A}(\lambda_k)\hat{\mathbf{u}},
\qquad
\mathcal{A}(\lambda_k)=\mathcal{J}_F-\lambda_k(\mathcal{D}-\mathcal{H}).
\]
For each fixed $k\ge2$, $\mathcal{A}(\lambda_k)$ is a real $2\times2$ matrix.
By the Routh--Hurwitz criterion for $2\times2$ systems, the mode $k$ is unstable if and only if
\begin{equation}
\label{eq:RH_2x2}
\operatorname{tr}\mathcal{A}(\lambda_k)>0
\qquad \text{or}\qquad
\det\mathcal{A}(\lambda_k)<0,
\end{equation}
with equalities corresponding to neutral stability boundaries.

A direct computation gives the trace formula \eqref{eq:twoway_trace_condition}, proving the trace criterion.
For the determinant, expanding $\det(\mathcal{A}(\lambda))$ as a polynomial in $\lambda$ yields \eqref{eq:det_quadratic_def}
with coefficients \eqref{eq:A1A2_def}. Since $\det(\mathcal{J}_F)>0$ by \eqref{eq:kinetic_stability_twoway}, the constant term is positive.
Thus $\det(\mathcal{A}(\lambda))<0$ for some $\lambda>0$ occurs exactly in the three subcases listed in Theorem~\ref{thm:twoway_criteria}:
\begin{enumerate}[label=(\roman*)]
\item If $A_1<0$, then $\det(\mathcal{A}(\lambda))\to -\infty$ as $\lambda\to\infty$, hence some large mode destabilizes.
\item If $A_1=0$, then $\det(\mathcal{A}(\lambda))=A_2\lambda+\det(\mathcal{J}_F)$, which is negative for some $\lambda>0$ iff $A_2<0$.
\item If $A_1>0$, then the upward-opening quadratic is negative for some $\lambda>0$ iff it has two real roots and both are positive.
Since the product of the roots is $\det(\mathcal{J}_F)/A_1>0$, the roots (when real) have the same sign; thus both are positive iff
their sum $-A_2/A_1$ is positive, i.e.\ $A_2<0$, and the discriminant is positive, i.e.\ $A_2^2-4A_1\det(\mathcal{J}_F)>0$.
\end{enumerate}
Combining these determinant facts with \eqref{eq:RH_2x2} completes the proof.
\end{proof}

\end{document}
